\subsection{群胚箭的缩并}

设 G 为群胚，F 为 G 的箭集的一个子集。设 H 为由 F 生成的 G 的正规子群胚。群胚 G/H 称为由 G 缩并箭集 F 而得的群胚，记作 G/F。若以 p 表示从 G 到 G/F 的典范态射，则有 \(p(o(f)) = p(t(f))\) 与 \(p(f) = e_{p(o(f))}\) （对每条箭 \(f \in F\) ）。

此态射满足下列泛性质。

命题 7. —— 对任何群胚态射 \(\varphi\colon G\to G'\) （它把属于 F 的每条箭都映为 \(G'\) 的单位元素），存在唯一的群胚态射 \(\varphi'\colon G/F\to G'\) ，使得 \(\varphi'\circ p=\varphi\) 。

\(\varphi\) 的核是 G 的正规子群胚，其箭集包含 F。按定义，H 是 G 的箭集包含 F 的最小正规子群胚；因此它含于 \(\varphi\) 的核。命题由 II, 第 170 页, 命题 3 得出。

以 \(\Gamma\) 表示 G 的子箭图，其顶点集为 Som(G)，箭集为 F。H 的轨道等同于箭图 \(\Gamma\) 的连通分支。按群胚对正规子群胚的商的定义，G/F 的顶点集等同于 \(\Gamma\) 的连通分支所成之集；换言之，它是 Som(G) 关于最细等价关系的商集，该关系使得 \(o(f)\) 与 \(t(f)\) 等价（对每条箭 \(f \in \mathrm{F}\) ）。

映射 p 通过取商诱导从 G 的轨道集到 G/F 的轨道集上的一个双射（II, 第 170 页, 注 1）。

本节的目的，是计算由 p 通过取迷向群而诱导的同态；据命题 2（II, 第 170 页），这等价于计算子群胚 H 的迷向群。

命题 8. —— 设 δ 为从自由群胚 \(\mathrm{Grp}(\Gamma)\) 到 G 中的典范态射，它开拓 Γ 到 G 中的典范单射态射。设 a 为 G 的顶点，A 为 a 在 G 中的轨道。对每个 b ∈ A，设 \(f_{ab}\) 为 G 中连接 a 与 b 的箭。迷向群 \(H_{a}\) 是 \(G_{a}\) 的正规子群，由元素 \(\operatorname{Int}(f_{ab})(\delta(c))\) 生成，其中 b 取遍集合 A，c 取遍 \(\mathrm{Grp}(\Gamma)_{b}\) 的元素。

设 x 与 y 为 G 中属于同一轨道的顶点，且 \(x \neq a\) ；我们再固定一条连接 x 与 y 的 G 的箭 \(f_{xy}\) 。对 G 的每个点 x，以 \(N_{x}\) 表示 \(G_{x}\) 的由元素 \(\operatorname{Int}(f_{xy})(\delta(c))\) 生成的正规子群，其中 y 取遍 x 在 G 中的轨道，c 取遍群 \(\operatorname{Grp}(\Gamma)_{y}\) 。对一切 \(f \in G_{xy}\) 以及每个 \(g \in N_{y}\) ，有 \(fgf^{-1} \in N_{x}\) 。设 \(x \in \operatorname{Som}(G)\) 。我们有 \(\delta(\operatorname{Grp}(\Gamma)_{x}) \subset N_{x}\) （按 \(N_{x}\) 的定义）。按群胚 H 的定义，\(\delta(f)\) 是 H 的箭——对 \(\operatorname{Grp}(\Gamma)\) 的每条箭 f 皆是；因而 \(\delta(\operatorname{Grp}(\Gamma)_{x})\) 含于 \(H_{x}\) 。由于 H 是 G 的正规子群胚，我们于是有

\[
\operatorname{Int} (f _ {x y}) (\delta (c)) \in \operatorname{Int} (f _ {x y}) (\mathrm{N} _ {y}) \subset \operatorname{Int} (f _ {x y}) (\mathrm{H} _ {y}) \subset \mathrm{H} _ {x}
\]

对一切 \(y \in \operatorname{Som}(\mathbf{G})\) ，便有 \(\mathrm{N}_x \subset \mathrm{H}_x\) 。

设 \(x\) 与 \(y\) 为 \(\Gamma\) 的顶点。令 \(\mathrm{N}_{x,y} = \varnothing\) （若 \(x\) 与 \(y\) 不属于 \(\Gamma\) 的同一连通分支）。否则，设 \(c\) 与 \(c'\) 为连接 \(x\) 与 \(y\) 的道路类（在图 \(\widetilde{\Gamma}\) ——与 \(\Gamma\) 相伴——中）。设 \(z\) 为 G 中属于 \(y\) 的轨道的顶点，\(\ell\) 为 \(\mathrm{Grp}(\Gamma)_z\) 的元素；在群 \(\mathrm{G}_x\) 中有等式：

\[
\begin{array}{r l} & {\delta (c) f _ {y z} \delta (\ell) f _ {y z} ^ {- 1} \delta (c ^ {\prime}) ^ {- 1}} \\ & {\qquad = \delta (c) f _ {y z} f _ {x z} ^ {- 1} \big (f _ {x z} \delta (\ell) f _ {x z} ^ {- 1} \big) f _ {x z} f _ {y z} ^ {- 1} \delta (c) ^ {- 1} \delta (c (c ^ {\prime}) ^ {- 1}).} \end{array}
\]

箭 \(\delta(c(c')^{-1})\) 是 \(x\) 处的一个圈（在箭图 \(\Gamma\) 中）的类，故属于 \(\mathrm{N}_x\) ，箭 \(f_{xz}\delta(\ell)f_{xz}^{-1}\) 亦然。由于 \(\delta(c)f_{yz}f_{xz}^{-1}\) 属于 \(\mathrm{G}_x\) ，且 \(\mathrm{N}_x\) 是 \(\mathrm{G}_x\) 的正规子群，故 \(\delta(c)f_{yz}\delta(\ell)f_{yz}^{-1}\delta(c')^{-1}\) 属于 \(\mathrm{N}_x\) 。从而 \(\delta(c)\mathrm{N}_y\subset \mathrm{N}_x\delta(c')\) 。由对称性，有 \(\delta(c)\mathrm{N}_y = \mathrm{N}_x\delta(c')\) 。这蕴含该集合不依赖于 \(c\) 与 \(c'\) 的选取；把它记作 \(\mathrm{N}_{x,y}\) 。

设 N 为 G 的子箭图，其顶点集为 Som(G)，且其从 x 到 y 的箭集等于 \(N_{x,y}\) （对 G 的每一对顶点 \((x,y)\) ）。它是 G 的正规子群胚，其箭集包含 F；因而 H 是 N 的子群胚。特别地，\(H_{a} \subset N_{a}\) ，由此得到等式。

推论 1. —— 再设对 F 的每条箭 f 有 \(o(f) = t(f)\) 。那么箭图 \(\Gamma\) 中每条道路的起点与靶相重合。设 a 为 G 的顶点，A 为其在 G 中的轨道。对 F 的每个起点 b 属于 A 的元素 c，固定一条连接 a 与 b 的 G 的箭 \(f_{c}\) ，并令 \(\kappa(c) = \operatorname{Int}(f_{c})(\delta(c)) = f_{c}\delta(c)f_{c}^{-1}\) 。

群 \(H_{a}\) 这时就是 \(G_{a}\) 的由箭 \(\kappa(c)\) 生成的正规子群。

由命题 8，诸箭 \(\kappa(c)\) 都含于群 \(H_a\) 。设 \(N_a\) 为 \(G_a\) 的包含它们的最小正规子群。按此命题，只需证明：对 G 的每个属于 A 的顶点 \(b\) 、每条箭 \(f_{ab}\) （它在 G 中连接 \(a\) 与 \(b\) ）以及每个元素 \(c\) （\(\mathrm{Grp}(\Gamma)_b\) 的），\(\operatorname{Int}(f_{ab})(\delta(c))\) 属于 \(N_a\) 。由于 F 的每个元素的起点与靶相等，顶点 \(b\) 处的圈（在图 \(\widetilde{\Gamma}\) 中）形如 \((b,(f_1,\varepsilon_1),b,\ldots,(f_n,\varepsilon_n),b)\) ，其中对每个 \(i\in\{1,\ldots,n\}\) ，\(f_i\) 是以 \(b\) 为起点的 F 的元素，且 \(\varepsilon_i\in\{-1,1\}\) 。我们有

\[
\operatorname{Int} (f _ {a b}) (\delta (c)) = f _ {a b} f _ {c} ^ {- 1} \kappa (f _ {1}) ^ {\varepsilon_ {1}} \dots \kappa (f _ {n}) ^ {\varepsilon_ {n}} f _ {c} f _ {a b} ^ {- 1}.
\]

由于 \(f_{ab}f_c^{-1} \in \mathrm{G}_a\) ，故 \(\operatorname{Int}(f_{ab})(\delta(c))\) 属于 \(\mathrm{N}_a\) 。

推论 2. —— 若与箭图 \(\Gamma\) 相伴的图是森林，则同态 \(p_a\) （从 \(\mathrm{G}_a\) 到 \((\mathrm{G} / \mathrm{F})_{p(a)}\) ）是双射，对 G 的每个顶点 a 皆然。

事实上，在此假设下，对 G 的每个顶点 b 有 \(\mathrm{Grp}(\Gamma)_{b}=\{e_{b}\}\) （II, 第 173 页, 推论 1）。于是由命题 8 推出，对 G 的每个顶点 a，群 \(H_{a}\) 退缩为单位元。

